% REVIEW DRAFT: not an input of the current published manuscript.
% This file is not proof certification; see geometry-audit-2026-11-21.json.

\section{Cuspidal orthogonality for the theta comparison space}
\label{sec:sr-theta-cuspidal-orthogonality}

Let \(\Sigma\) be the span of the large polynomial theta periods
and the small polynomial theta periods in the actual cutoff,
for one fixed triple. By the actual period factorization each
local module of \(\Sigma\) is a quotient of a finite sum of
copies of the common matching module. The target periods are
square integrable. We prove the orthogonality assertion needed
after a difference in \(\Sigma\) has been shown to be cuspidal.

We use the following precisely applicable published rapid-decrease
theorem. Miller--Schmid, Theorems~2.23--2.24, prove that a smooth
cuspidal automorphic form of moderate growth and finite
generalized infinitesimal character on an arithmetic quotient
of a real reductive group with compact center has uniform
moderate growth and decays faster than every inverse norm
power on each generalized Siegel set. Their statement allows
smooth vectors without compact finiteness and includes the
corresponding seminorm bounds.
\footnote{S.~Miller and W.~Schmid, \emph{On the rapid decay of
cuspidal automorphic forms}, Adv.~Math.~231 (2012), 940--964,
Theorems~2.23--2.24:
\url{https://arxiv.org/pdf/1106.2149}.}

Here the real group is
\((\operatorname{Res}_{F/\mathbb Q}U(W))(\mathbb R)\);
its center is a product of circles, hence compact.
Fix the finite level of \(F\in\Sigma\).
The adelic quotient has finitely many arithmetic real
components at that level. On each, unipotent strong
approximation identifies the adelic zero constant terms
with the arithmetic real zero constant terms.
Thus the form is cuspidal in the theorem's sense on each
component. The actual real common matching module gives
its finite generalized infinitesimal character.
The proved theta bounds give moderate growth of its value;
no uniform derivative growth assumption is added.
The theorem supplies that conclusion and rapid decrease.
Its matrix norm and our height are polynomially comparable
on every fixed finite component, so its arbitrary inverse
norm powers are precisely the required arbitrary inverse
height powers. It also applies to fixed right translates,
which retain moderate growth, the central annihilator
and all the zero constant terms.
This verifies the rapid-decrease input for every cuspidal
\(F\in\Sigma\), rather than postulating it for the difference.

\begin{srlemma}[The actual back-lifts are constant]
\label{lem:sr-theta-backlift-constant-all-signatures}
Let \(F\in\Sigma\) be cuspidal. For a Schwartz vector \(\Psi\)
in either upper oscillator model put
\[
 \mathcal T_\Psi F(g)
     =\int_{[U(W)]}F(h)\,
                    \overline{\theta(\Psi;g,h)}\,dh .
 \tag{O1}\label{eq:sr-theta-backlift-all-signatures}
\]
This converges, is continuous and automorphic in \(g\), and
is a constant function on its first unitary group.
\end{srlemma}

\begin{proof}
At a fixed \(g\), the kernel as a function of \(h\) is in \(L^2[U(W)]\).
Indeed the all-rank lattice estimate with the first argument
fixed bounds it on a target Siegel set by
\(C_g\exp\{Dn\sum t_j/2\}\).
Its squared integral has initial coefficient sums
\(Dk(n-d+k)<0\), because the target Witt index is at most
\(Y=p+q\), whereas
\(n+Y<d\) for both \(n=a\) and \(n=a'\).
The constants are uniform when \(g\) varies in a compact set.
Cauchy--Schwarz with \(F\in L^2\) proves convergence and local
uniform domination. Schwartz continuity gives continuity.
Rational invariance of theta gives automorphy.

We now apply the good split matrix-functional calculation already
proved in Section~\ref{sec:sr-compact-global-comparison}.
Here is the needed extension of that calculation beyond the
compact base. Fix a good fully split place other than the cutoff
place and freeze all the outside tensor factors.
For a back-lift to the rank-\(a\) first space, parameterize the
local common class of \(F\) by the \emph{onto primed} quotient
on \(a'\)-row Schwartz matrices. For a back-lift to the
rank-\(a'\) space, parameterize it by the \emph{onto unprimed}
quotient on \(a\)-row matrices. Thus the two local test
variables always have exactly \(a+a'=d\) rows.
This applies to a sum of the two images as well, since both
period maps have already descended to that same common module.

The resulting bilinear functional is equivariant for the common
right \(\mathrm{GL}_d\) action. It is an ordinary algebraic
functional on the finite local Schwartz spaces: its values are
the convergent \(L^2\) pairings just proved.
The rank filtration, stabilizer modular mismatch, and full-rank
fiber calculation from that section therefore apply verbatim.
They prove invariance of (O1) under the first local unitary group
at every such good split place. Compactness of the entire
automorphic quotient was not used in that local calculation;
only convergence of the displayed pairing was needed.

Finally the already proved good-place density argument applies.
Strong approximation for the special unitary group away from one
noncompact split finite place reduces a continuous automorphic
function invariant at these places to the norm-one determinant
torus. That adelic torus quotient is compact. A nontrivial
character trivial at all good split places would give a
nontrivial unitary Hecke character whose Euler factors at all
degree-one rational primes outside a finite set are those of
the zeta function. Since \(E/\mathbb Q\) is Galois, the remaining
prime ideals have rational residue degree at least two.
The logarithm of their Euler-product contribution is bounded
as \(s\downarrow1\), contradicting entire nontrivial Hecke
\(L\)-functions and their nonvanishing at \(1\).
Thus these local groups are dense on the torus quotient also.
The back-lift is constant, as claimed.
\end{proof}

\begin{srproposition}[Orthogonality when the first space is isotropic]
\label{prop:sr-theta-orthogonal-isotropic}
Let \(F\in\Sigma\) be cuspidal.
If one of the first Hermitian spaces is rationally isotropic,
then \(F\) is orthogonal to every polynomial theta period from
that space, with the actual cutoff on the small side.
\end{srproposition}

\begin{proof}
It suffices to test vectors in the actual full-rank cutoff.
For the small period this is its given domain.
For a large period replace its finite factor at the good cutoff
place by a full-rank representative of the same matching class.
The exact full-rank coinvariant isomorphism supplies such a
representative; the proved period factorization says that
the large period is unchanged. Finite sums are treated by
the same tensor presentation argument as before.

For this cutoff vector \(\Psi\), full-rank decay gives, for every
\(N\),
\[
 |\theta(\Psi;g,h)|
       \le C_N\,\operatorname{qht}_X(g)^{-N}
                         \operatorname{ht}_W(h)^{b_N}.
\]
Choose the target Siegel fundamental set. On it its ordinary
height is bounded by a fixed power of quotient height:
the upper chamber height bound and the rational-translate
product-formula lower bound proved in the decay section give
this comparison. Rapid decrease of \(F\) now gives
\[
 |\mathcal T_\Psi F(g)|
       \le C'_N\operatorname{qht}_X(g)^{-N}.
 \tag{O2}\label{eq:sr-theta-backlift-decay}
\]
The same estimate justifies Fubini with the first quotient
integral, since that quotient has finite Haar volume.

There is an unbounded quotient-height sequence in the isotropic
first group. In its rational Witt basis take a chamber parameter
tending to infinity. A nonzero entry in a negative row of a
rational inverse matrix, together with the product formula,
gives the rational-translate lower height bound; taking the
infimum over rational translates still tends to infinity.
By the preceding lemma the left side of (O2) is a constant.
Letting this sequence tend to infinity shows that constant is zero.
Fubini consequently gives
\[
 \langle I_X(\Psi),F\rangle
        =\int_{[U(X)]}\overline{\mathcal T_\Psi F(g)}\,dg=0.
\]
This proves the assertion for all the periods under consideration.
\end{proof}

\begin{srlemma}[Cuspidal pullback of the doubled Eisenstein kernel]
\label{lem:sr-theta-cuspidal-doubling-pullback}
Let \(X\) be definite and let \(F\) be a rapidly decreasing
cuspidal form on \(U(W)\).
For an actual doubled flat Schwartz section, integrate its
Eisenstein kernel, with the specified splitting determinant
factor, against \(F\) in its second \(U(W)\) variable.
The resulting meromorphic family in the first variable is
cuspidal wherever it is regular.
\end{srlemma}

\begin{proof}
Write \(G=U(W)\), \(H=U(W\oplus(-W))\), and \(P\subset H\)
for the diagonal maximal-isotropic stabilizer.
The doubled flat section is right compact finite, smooth at
infinity, and locally constant of finite level at finite places.
Only finitely many finite factors are nonstandard.
Its inducing determinant character has absolute value one.
If the prescribed doubling normalization multiplies the
restricted Eisenstein series by separate determinant
characters of the two copies, absorb the second character
into \(F\) and restore the first after integration.
A unitary determinant twist preserves rapid decrease and
every zero unipotent constant term. Thus it suffices to
prove the statement for the raw restricted section;
the total determinant factors are retained in
\(\vartheta\) and \(k_s\) below.

First work with sufficiently large real part of \(s\).
The orbit argument can be checked without an unspecified
kernel identity. Let \(L\) be a maximal isotropic subspace of
\(W\oplus(-W)\), and put
\[
 L_2=L\cap(0\oplus W),\qquad Y_2=\operatorname{pr}_2 L .
\]
Isotropy implies \(Y_2\subset L_2^\perp\).
If \(L_2\ne0\), the second-copy unipotent radical
\(N_{L_2}\) fixes \(L_2\) and acts trivially on
\(L_2^\perp/L_2\). Consequently
\((1,n)L=L\) for every \(n\in N_{L_2}\):
the change of any second projection lies in \(L_2\),
which is itself contained in \(L\).
The induced action on \(L\) is unipotent, with determinant one.
The inducing character and its half modulus therefore are
trivial on this stabilizer. Unfolding this orbit integrates
\(F\) first along
\(N_{L_2}(F)\backslash N_{L_2}(\mathbb A_F)\);
cuspidality makes it zero.
For a maximal doubled isotropic subspace, the dimensions of
its two intersections agree: if their dimensions are \(r_1,r_2\),
its two projections have dimensions \(d-r_1,d-r_2\), and
\(Y_i\subset L_i^\perp\) gives both inequalities
\(r_1\ge r_2\) and \(r_2\ge r_1\).
Thus every nonopen orbit has \(L_2\ne0\) and is covered.

If both intersections vanish, the two projections identify
\(L\) with the graph of a rational isometry of \(W\).
Witt extension gives one rational open orbit, with diagonal
\(G(F)\) stabilizer. On this orbit the actual section satisfies
\[
 f_s(\delta\iota(h,h'))
       =\vartheta(h)\,
                      f_s(\delta\iota(1,h^{-1}h')),
\]
where \(\vartheta\) is the specified determinant splitting
character on \(G(\mathbb A_F)\). Its restriction to every
unipotent subgroup is one. Indeed simultaneous left
multiplication by \((h,h)\) stabilizes the diagonal
isotropic subspace and acts on it by \(h\); apply the
induction equivariance. The adelic norm of its determinant
is one, so no \(s\)-dependent norm factor remains.
The unfolded pullback is therefore exactly
\[
       \vartheta(h)\int_{G(\mathbb A_F)}
                                k_s(g)F(hg)\,dg,\qquad
       k_s(g)=f_s(\delta\iota(1,g)).                \tag{O4}
 \label{eq:sr-theta-cusp-open-orbit}
\]

All the unfoldings just made are absolute in a sufficiently
positive half-plane. One may check this directly on the
Cartan decomposition already used for the theta estimates:
the diagonal-isotropic Iwasawa determinant of
\(\delta\iota(1,g)\) decreases by a fixed positive power of
each growing Cartan coordinate. Thus \(k_s(g)\) is bounded
by the corresponding negative powers with exponents tending
to infinity with \(\operatorname{Re}s\).
The Cartan Haar density has only fixed positive root powers.
At a good finite place its resulting absolute integral is
\(1+O(q_v^{-\operatorname{Re}s+c})\), with \(c\) depending
only on the ranks; finitely many bad places have convergent
geometric-series bounds for large \(\operatorname{Re}s\).
The infinite Cartan integrals converge by the same comparison.
Hence \(k_s\in L^1(G(\mathbb A_F))\) there.
Equivalently, the absolutely convergent positive majorant
for the Eisenstein series has polynomial growth on the
target Siegel set, which is integrable against the rapid
cusp form. This also justifies orbit unfolding and its
compact unipotent averages. This is also the ordinary
doubling unfolding, for arbitrary unitary signatures:
Eischen's Theorem~4.3.4 and its orbit proof in
\S4.3.4 give this same graph-orbit calculation.
\footnote{E.~Eischen, \emph{Automorphic Forms on Unitary Groups},
AWS notes, \S4.3.4:
\url{https://swc-math.github.io/aws/2022/2022EischenNotes.pdf}.}
The first-variable constant term of (O4) is zero:
move that compact average inside its absolutely convergent
integral, use \(\vartheta(n)=1\), and apply the zero
constant term of \(F\) at \(hg\).
This proves cuspidality in the initial half-plane.

We next justify the continuation under these integrals.
This requires more than pointwise meromorphic continuation.
Bernstein--Lapid's corollary at the end of \S8, p.~226,
states that locally a scalar polynomial clears the poles
and the resulting Eisenstein family is holomorphic in one
Fr\'echet space of smooth functions with a fixed uniform
moderate-growth exponent. It applies to this compact-finite
flat section and its determinant induction.
\footnote{J.~Bernstein and E.~Lapid,
\emph{On the meromorphic continuation of Eisenstein series},
JAMS~37 (2024), 187--234, Theorem~2.3 and the corollary following
its proof, author's offprint:
\url{https://www.math.tau.ac.il/~bernstei/Publication_list/publication_texts/2023-10-15-Eisenstein-series-JAMS.pdf}.}
If \(h\) lies in a fixed compact set and \(h'\) lies in a
target Siegel set, restriction of its defining seminorm
therefore bounds the pole-cleared kernel by
\(C\,\operatorname{ht}(h')^R\), with one \(R\), uniformly
on a compact parameter set. Rapid decrease of \(F(h')\)
makes
\[
 K\longmapsto
       \int_{[G]}K(\iota(h,h'))F(h')\,dh'
\]
a continuous linear functional on that Fr\'echet space,
uniformly for such \(h\). Averaging the first variable
over a rational unipotent quotient is another such continuous
functional: that quotient has compact representatives.
Applying these functionals to the holomorphic pole-cleared
family proves their holomorphy, hence meromorphy before
clearing the poles. Their zero in the initial half-plane
continues to zero everywhere.

At a Siegel--Weil special value the definite, hence anisotropic,
identity used below asserts regularity of the actual section's
Eisenstein series, not a renormalized residue. This is enough
also in the Fr\'echet sense: every coefficient in its possible
finite principal Laurent part vanishes at every group point,
so all those coefficient functions are zero.
Thus evaluating the integral and the zero first constant term
at that value is legitimate. The same argument applies at
any other regular value. No isotropic Siegel--Weil assertion
enters this continuation argument.
\end{proof}

\begin{srproposition}[Orthogonality for a definite first space]
\label{prop:sr-theta-orthogonal-definite}
Suppose the target \(W\) is isotropic, one first space \(X\)
is definite, and an adapted opposite-rank tower step is available.
Suppose the lower period on that unchanged definite first space
is nonzero. Then a cuspidal \(F\in\Sigma\) is orthogonal to every
polynomial theta period from \(X\).
For the small definite first space its tested domain may be
restricted to the actual cutoff.
\end{srproposition}

\begin{proof}
Assume instead that
\(\langle I_X(\Psi_2),F\rangle\ne0\).
The back-lift \(\mathcal T_{\Psi_2}F\) is a constant \(b\).
Since \([U(X)]\) is compact, ordinary Fubini gives
\[
         \langle F,I_X(\Psi_2)\rangle
                      =\operatorname{vol}[U(X)]\,b;
\]
in particular \(b\ne0\).

We check the compact-finiteness required by the classical
identity and continuation, rather than assuming it for an
arbitrary box vector. At infinity the Cartan on the doubled
target \(W\oplus(-W)\) is
\(\operatorname{diag}(\theta_W,-\theta_W)\).
Its positive oscillator polarization is the direct sum of
the polarization for \(X\otimes W\) and its complex conjugate.
Thus \(\Psi_1\boxtimes\overline{\Psi_2}\), for polynomial
Gaussian \(\Psi_1,\Psi_2\), is a polynomial times this doubled
positive Gaussian. The doubled maximal compact fixes its
Gaussian line and preserves each bounded polynomial degree,
so the box is compact finite for the full doubled group.
At a finite place smoothness of the Weil representation
gives an open stabilizer of its Schwartz--Bruhat vector;
its orbit under a compact subgroup is finite.
Almost all factors are the unramified vacuum and are fixed.
These facts prove compact-finiteness of the global section.
For definite \(X\), its own Cartan is \(+1\) or \(-1\)
according to its sign, so the stated positive polarizations
are the actual polynomial spaces, with no varying choice.

Apply the proved anisotropic Siegel--Weil identity to this
\emph{definite} \(X\) and this actual doubled vector.
With the already checked Haar normalization its special value is
\[
 E_{\Psi_1,\Psi_2}(h,h')
   =c_{\rm SW}\operatorname{vol}[U(X)]^{-1}
       \int_{[U(X)]}\theta(\Psi_1;g,h)
                    \overline{\theta(\Psi_2;g,h')}\,dg ,
 \tag{O3}\label{eq:sr-theta-anisotropic-doubled-kernel}
\]
including the prescribed splitting determinant factor.
Here \(c_{\rm SW}\) is the nonzero scalar \(1\) or \(2\)
already calculated. Compactness of \([U(X)]\), its uniform
polynomial theta bound, and rapid decrease of \(F\) justify
integrating (O3) against \(F(h')\) and interchanging integrals.
The result is
\[
    c_{\rm SW}\operatorname{vol}[U(X)]^{-1}\,
                                      b\,I_X(\Psi_1)(h).
\]
The preceding pullback lemma says this function is cuspidal.
Its scalar is nonzero, so \emph{every} polynomial period
\(I_X(\Psi_1)\) is cuspidal.

This contradicts the already proved zero-term formula in the
opposite-rank tower. In that step \(X\) is unchanged, the target
loses one hyperbolic plane, and the zero constant-term component
is a fixed nonzero constant times the lower period of the
actual zero evaluation. That evaluation is onto on polynomial
Gaussians and finite Schwartz vectors. On the small side
a lower cutoff vector has an upper cutoff preimage: in its
split Fourier coordinates retain the full-rank lower block.
Multiply the fixed zero column by a compact open bump of
value one at zero, and the integrated column by a compact
open bump of integral one. The zero-evaluation formula
\(e_0\phi(C)=c_v\int\phi(0,C,z)\,dz\), with its fixed
nonzero \(c_v\), then gives the desired lower vector after
one scalar rescaling.
Thus lower nonvanishing gives a test vector with a nonzero
upper constant-term component. An upper period cannot then
be cuspidal. This contradiction proves orthogonality.
\end{proof}

\begin{srcorollary}[A cuspidal comparison difference vanishes]
\label{cor:sr-theta-cuspidal-difference-zero}
Suppose the lower nonvanishing required by the definite
orthogonality proposition is available.
If \(F\in\Sigma\) is cuspidal, then \(F=0\).
\end{srcorollary}

\begin{proof}
In a noncompact target at least one of the signature parameters
\(p,q\) is positive. A first space with positive Witt index
is covered by the isotropic orthogonality proposition.
A definite first space is covered by the other proposition,
using the opposite-rank step. The needed rational isotropy
or adapted step is the one in the tower construction.
Hence \(F\) is orthogonal to both theta-image summands which
span \(\Sigma\). In particular \(\langle F,F\rangle=0\).
It is zero as an \(L^2\) function and, by continuity, pointwise.
The compact target is the already proved compact-base
comparison and is not an application of the present cusp
argument.
\end{proof}

Only the definite case used Siegel--Weil here, and it used its
proved anisotropic form with the actual box section and Haar
normalization. No regularized isotropic Siegel--Weil identity
has been inserted into the argument that is meant to prove
the global theta comparison.
